% ============================================================================
%  9. Decisiones de diseño y discrepancias resueltas
% ============================================================================
\section{Decisiones de diseño y discrepancias resueltas}
\label{sec:decisiones}

Al contrastar la guía del taller con las hojas de datos se encontraron cinco
puntos que, sin resolverse, habrían afectado el funcionamiento o la evaluación
(tabla~\ref{tab:decisiones}).

\begin{table}[H]
  \centering
  \caption{Hallazgos y decisiones de diseño.}
  \label{tab:decisiones}
  \small
  \begin{tabularx}{\textwidth}{@{}c L L L@{}}
    \toprule
    \textbf{\#} & \textbf{Hallazgo} & \textbf{Impacto} & \textbf{Decisión} \\
    \midrule
    1 & \SI{1.2}{\kilo\hertz} con \SI{20}{\mega\hertz} y prescaler 1:16 exige $PR2 = 259.4$, fuera del rango de 8~bits. &
        Imposible de configurar. &
        $PR2 = 249 \rightarrow \SI{1250}{\hertz}$ (error 4.2\,\%), con escala exacta de 1000 pasos (sección~\ref{sec:calculos}). \\
    2 & La lista de materiales indica cristal de \SI{4}{\mega\hertz}; el código usa \SI{20}{\mega\hertz} (FOSC = HS). &
        Con \SI{4}{\mega\hertz} el PWM baja a \SI{250}{\hertz} y los retardos se alargan 5 veces. &
        Directiva \texttt{CRISTAL\_20MHZ} que ajusta FOSC, prescaler y SPBRG para ambos casos. \\
    3 & La rúbrica evalúa el pulso de activación en nSLEEP, pero el código base no maneja ese pin. &
        El driver depende de un pull-up del módulo o queda deshabilitado. &
        nSLEEP en RD2, en 0 durante la configuración y flanco $0 \rightarrow 1$ con \texttt{DRV8833\_Wake()}. \\
    4 & La rúbrica evalúa giro, pero la secuencia base solo tiene avance, reversa y freno. &
        Criterio de demostración incompleto. &
        Giro a la izquierda y a la derecha sobre su eje agregados a la secuencia. \\
    5 & Capacitores de cristal de \SI{44}{\pico\farad} en la lista de materiales. &
        Fuera del rango recomendado de 15 a \SI{33}{\pico\farad}; arranque del oscilador poco confiable. &
        Usar \SI{22}{\pico\farad}, como en el esquemático. \\
    \bottomrule
  \end{tabularx}
\end{table}

También se definió el sentido de giro según la convención física (izquierda =
rotación antihoraria vista desde arriba), de modo que el comando y el
movimiento observado coincidan durante la demostración.
